\section{Plan de Acción Preliminar}
\apendiceindice{D}{Plan de Acción Preliminar}
Las fechas son una propuesta sujeta al calendario escolar y a la
autorización del plantel; coinciden con el calendario propuesto en el bosquejo del
método.
\begin{enumerate}
	\item \textbf{Diagnóstico (18 de noviembre al 4 de diciembre de 2026):}
	      solicitar las autorizaciones, revisar la legislación de protección de
	      datos (Apéndice~F), probar los instrumentos y aplicar la encuesta, las
	      entrevistas y la primera serie de observación de los Apéndices A a C
	      para obtener la línea base del acceso piloto.
	\item \textbf{Diseño del sistema, hardware y firmware:} definir la
	      arquitectura y el modelo entidad-relación, integrar el ESP32-CAM con el
	      mecanismo del torniquete y programar la lectura del código y la consulta
	      al servidor.
	\item \textbf{Desarrollo del servidor:} implementar la generación y
	      validación de QR dinámicos firmados con HMAC, la base de datos de ingresos y
	      el panel web con los indicadores definidos.
	\item \textbf{Pruebas de banco (30 de noviembre al 4 de diciembre de
	      2026):} aplicar el protocolo del Apéndice~E y verificar que el prototipo
	      cumple los criterios de aceptación antes de instalarlo.
	\item \textbf{Piloto controlado (7 al 15 de diciembre de 2026):} instalar el
	      prototipo, dar de alta a los estudiantes que firmen el consentimiento del
	      Apéndice~G, dejar dos días de adaptación y levantar la segunda serie de
	      observación durante cinco días hábiles. El código QR sustituye la
	      revisión visual; la revisión manual queda solo para quien no pueda
	      mostrar su código o no participe en el piloto.
	\item \textbf{Evaluación (16 de diciembre de 2026 al 29 de enero de 2027):}
	      comparar la segunda serie con la línea base, analizar la encuesta y las
	      entrevistas, documentar los ajustes necesarios y presentar los
	      resultados.
\end{enumerate}
